% =====================================================================
% Premessa: come leggere il manuale (legenda dei riquadri e fonti)
% =====================================================================
\chapter*{Come leggere questo manuale}
\addcontentsline{toc}{chapter}{Come leggere questo manuale}

Questo manuale è scritto da studenti, per studenti: \textbf{non è materiale ufficiale del corso} e non
sostituisce le slide, i testi normativi e le indicazioni del docente. Nasce dalle registrazioni delle
lezioni (trascritte automaticamente e poi revisionate), integrate con slide e materiale dei laboratori.

\section*{Le parti}
\begin{description}[leftmargin=2em,style=nextline]
  \item[Diario delle lezioni] una voce per lezione, in ordine cronologico: sintesi, avvisi pratici,
    contenuti in breve e punti segnalati per l'esame. Serve a seguire l'avanzamento del corso.
  \item[Manuale del corso] un capitolo per norma (GDPR, data breach, ISO/IEC~27001, CRA, NIS2, DORA,
    AI Act, CSA), che unisce lezioni e slide in un'esposizione ordinata. Rimanda alle lezioni del diario.
  \item[Laboratori] i casi pratici del venerdì: contesto, consegna, metodo e soluzione ragionata.
  \item[Schemi di ripasso] tabella delle scadenze, mappe concettuali, schede riassuntive.
  \item[Quiz ed esame] formato dell'esame, tema precedente, simulazioni e quiz per argomento con
    soluzioni e rimando alla lezione.
  \item[Appendice] glossario e acronimi.
\end{description}

\section*{Legenda dei riquadri}

\begin{esame}
Punto che il docente ha indicato come rilevante per l'esame.
\end{esame}

\begin{integrazione}
Contenuto aggiunto dai redattori: \textbf{non detto a lezione e non presente nelle slide}
(spiegazioni, esempi, riferimenti normativi in più). Utile, ma da trattare con cautela all'esame.
\end{integrazione}

\begin{verifica}
Informazione non confermata: passaggio poco chiaro nella registrazione, dato ricostruito o calcolato.
\end{verifica}

\begin{nota}
Nota generica (fonti usate, avvertenze, collegamenti).
\end{nota}

\begin{soluzione}
Soluzione di quiz, simulazioni e laboratori.
\end{soluzione}

\dacompletare{Parte ancora da scrivere (per esempio una norma non ancora trattata a lezione).}

\section*{Fonti}
Registrazioni delle lezioni (Prof.~Giuseppe Vaciago; laboratori con Gianluca Gilardi), slide e materiale
pubblicati sulla piattaforma del corso, testi normativi ufficiali (EUR-Lex). Le slide e le
registrazioni non sono incluse nel repository.
